% ============================================================
% NIH/DOE SciENcv — Biographical Sketch Common Form + NIH Supplement
% OMB No. 3145-0279 / OMB No. 0925-0001 & 0925-0002
%
% Jiachen (Amber) Liu — ORCID 0000-0003-2317-1956
% Prepared for: DE-FOA-0003612 (DISCOVER), Topic 20, Focus Area A
% "Adversarial Robustness Evaluation and Defenses for Science
%  Foundation Models and Agentic Workflows"
% PI: Prof. Atul Prakash, University of Michigan
% ============================================================
\documentclass[10pt]{article}

\usepackage[top=0.70in, bottom=0.70in, left=0.75in, right=0.75in]{geometry}
\usepackage{array, booktabs, tabularx, multirow}
\usepackage{enumitem}
\usepackage{hyperref}
\usepackage{xcolor}
\usepackage{fancyhdr}
\usepackage{parskip}
\usepackage{titlesec}
\usepackage[T1]{fontenc}
\usepackage{lmodern}
\usepackage{microtype}

\hypersetup{colorlinks=true, urlcolor=blue, linkcolor=blue}

% ── Section style ────────────────────────────────────────────
\titleformat{\section}[block]
  {\normalsize\bfseries}{}{0em}{}[\vspace{-4pt}\rule{\linewidth}{0.6pt}]
\titlespacing{\section}{0pt}{8pt}{4pt}

\titleformat{\subsection}[runin]{\small\bfseries\itshape}{}{0em}{}[.\enspace]
\titlespacing{\subsection}{0pt}{5pt}{4pt}

% ── Header / footer ─────────────────────────────────────────
\pagestyle{fancy}
\fancyhf{}
\fancyhead[L]{\small\textbf{Biographical Sketch Common Form}\enspace OMB No. 3145-0279 (Exp. 10/31/2026)}
\fancyhead[R]{\small Liu, Jiachen}
\fancyfoot[C]{\small\thepage}
\renewcommand{\headrulewidth}{0.4pt}

\setlength{\parindent}{0pt}
\setlength{\parskip}{3pt}
\newcolumntype{L}[1]{>{\raggedright\arraybackslash}p{#1}}

\begin{document}
\thispagestyle{fancy}

% ============================================================
%  PART 1 — BIOGRAPHICAL SKETCH COMMON FORM
% ============================================================

{\footnotesize OMB No.~3145-0279 \quad Expiration Date 10/31/2026}

\begin{center}
  {\large\bfseries BIOGRAPHICAL SKETCH COMMON FORM}\\[2pt]
  {\small Prepared for: DE-FOA-0003612 (DISCOVER), Topic 20, Focus Area A}
\end{center}
\vspace{-4pt}

% ── Identifying Information ──────────────────────────────────
\section*{Identifying Information}

\begin{tabular}{@{}L{5.2cm} L{5.0cm} L{4.8cm}@{}}
\textbf{*Name} (Last, First, Middle)  & \textbf{*Persistent Identifier (ORCID)} & \textbf{*Position Title} \\
Liu, Jiachen (Amber)                  & 0000-0003-2317-1956                     & Co-Founder \& CTO \\[6pt]
\multicolumn{3}{@{}l}{\textbf{*Organization and Location}} \\
\multicolumn{3}{@{}l}{Orchestra Research Institute, Inc.\quad|\quad 1 Broadway, 7th Floor, Cambridge, MA 02142, USA}
\end{tabular}

% ── Professional Preparation ─────────────────────────────────
\section*{Professional Preparation}
\textit{Reverse chronological order by start date. Includes all postdoctoral and fellowship training.}

\vspace{4pt}
\begin{tabular}{@{}L{4.4cm} L{3.2cm} L{4.2cm} L{1.3cm} L{2.0cm}@{}}
\toprule
\textbf{Organization} & \textbf{Location} & \textbf{Degree / Training} & \textbf{Start} & \textbf{End} \\
\midrule
University of Michigan  & Ann Arbor, MI, USA  & Ph.D., Computer Science        & 2020 & 05/2025 \\
\multicolumn{5}{@{}l}{\quad\small\itshape Advisor: Prof. Mosharaf Chowdhury (SymbioticLab) — ML Systems, LLM Serving, Federated Learning} \\[3pt]
MIT CSAIL (visiting)    & Cambridge, MA, USA  & Visiting Researcher             & 2019 & 01/2020 \\
\multicolumn{5}{@{}l}{\quad\small\itshape Database Systems / High-Dimensional Vector Search. Advisor: Prof. Samuel Madden} \\[3pt]
University of Michigan  & Ann Arbor, MI, USA  & B.S.E., Data Science            & 2018 & 05/2020 \\
\multicolumn{5}{@{}l}{\quad\small\itshape Minor in Mathematics} \\[3pt]
Shanghai Jiao Tong Univ.& Shanghai, China     & B.E., Elec.\ \& Comp.\ Eng.    & 2016 & 08/2020 \\
\bottomrule
\end{tabular}

% ── Appointments and Positions ───────────────────────────────
\section*{Appointments and Positions}
\textit{Reverse chronological order. All domestic and foreign professional appointments.}

\vspace{4pt}
\begin{tabular}{@{}L{1.2cm} L{1.3cm} L{5.0cm} L{4.5cm} L{2.5cm}@{}}
\toprule
\textbf{Start} & \textbf{End} & \textbf{Title} & \textbf{Organization} & \textbf{Location} \\
\midrule
2026 & Present & Co-Founder \& Chief Technology Officer & Orchestra Research Institute, Inc. & Cambridge, MA, USA \\
\multicolumn{5}{@{}l}{\quad\small\itshape AI-native research platform for agentic scientific workflows.} \\[3pt]
2025 & 2026    & Research Scientist & Meta Superintelligence Lab & Menlo Park, CA, USA \\
\multicolumn{5}{@{}l}{\quad\small\itshape LLM pre- and post-training systems (Llama).} \\[3pt]
2024 & 2024    & Research Scientist (Part-time) & Meta AI \& Systems Co-Design & Menlo Park, CA, USA \\
\multicolumn{5}{@{}l}{\quad\small\itshape Cross-regional Llama pre-training systems.} \\[3pt]
2024 & 2024    & Teaching Assistant & Univ.\ of Michigan, CSE & Ann Arbor, MI, USA \\
\multicolumn{5}{@{}l}{\quad\small\itshape EECS 598: Systems for Generative AI.} \\[3pt]
2022 & 2022    & Ph.D.\ Intern & Apple & New York, NY, USA \\
\multicolumn{5}{@{}l}{\quad\small\itshape Private Machine Learning Framework.} \\[3pt]
2019 & 2020    & Research Intern & MIT CSAIL Database Group & Cambridge, MA, USA \\
\multicolumn{5}{@{}l}{\quad\small\itshape High-dimensional vector database / ANN search.} \\[3pt]
2020 & 2025    & Graduate Research Assistant & Univ.\ of Michigan, SymbioticLab & Ann Arbor, MI, USA \\
\bottomrule
\end{tabular}

% ── Products ─────────────────────────────────────────────────
\section*{Products}

\subsection*{(i) Five products most closely related to the proposed project}

{\small
\begin{enumerate}[noitemsep, topsep=2pt, leftmargin=1.5em, label=\arabic*.]

  \item Kon P.*, \textbf{Liu J.*}, Zhu X., Ding Q., Peng J., Xing J., Huang Y., Qiu Y., Srinivasa J., Lee M.,
        Chowdhury M., Zaharia M., Chen A.
        \textit{EXP-Bench: Can AI Conduct AI Research Experiments?}
        ICLR 2026. \url{https://arxiv.org/abs/2505.24785}

  \item Kon P.*, \textbf{Liu J.*}, Ding Q., Qiu Y., Yang Z., Huang Y., Srinivasa J., Lee M., Chowdhury M., Chen A.
        \textit{Curie: Toward Rigorous and Automated Scientific Experimentation with AI Agents.}
        arXiv, 2025. \url{https://arxiv.org/abs/2502.16069}

  \item \textbf{Liu J.*}, Harmon M.*, Zhang Z.
        \textit{Sci-Reasoning: A Dataset Decoding AI Innovation Patterns.}
        arXiv, 2026. \url{http://arxiv.org/abs/2601.04577}

  \item Jabbour S., Chang T., Das Antar A., Peper J., Jang I., \textbf{Liu J.}, Chung J., He S., Wellman M.,
        Goodman B., Bondi-Kelly E., Samy K., Mihalcea R., Chowdhury M., Jurgens D., Wang L.
        \textit{Evaluation Framework for AI Systems in ``the Wild''.}
        arXiv, 2025. \url{https://arxiv.org/abs/2504.16778}

  \item Chung J., \textbf{Liu J.}, Ma J., Wu R., Kweon O.J., Xia Y., Wu Z., Chowdhury M.
        \textit{The ML.ENERGY Benchmark: Toward Automated Inference Energy Measurement and Optimization.}
        NeurIPS 2025 \textbf{(Spotlight)}. \url{https://arxiv.org/abs/2505.06371}

\end{enumerate}
}

\subsection*{(ii) Five other significant products}

{\small
\begin{enumerate}[noitemsep, topsep=2pt, leftmargin=1.5em, label=\arabic*., start=6]

  \item \textbf{Liu J.}, Wu Z., Chung J., Lai F., Lee M., Chowdhury M.
        \textit{Andes: Defining and Enhancing Quality-of-Experience in LLM-Based Text Streaming Services.}
        arXiv, 2024. \url{https://arxiv.org/abs/2404.16283}

  \item \textbf{Liu J.}, Lai F., Ding D., Zhang Y., Chowdhury M.
        \textit{Venn: Resource Management for Collaborative Learning Jobs.}
        MLSys 2025. \url{https://arxiv.org/abs/2312.08298}

  \item \textbf{Liu J.}, Lai F., Dai Y., Akella A., Madhyastha H., Chowdhury M.
        \textit{Auxo: Efficient Federated Learning via Scalable Cohort Identification.}
        SoCC 2023. \url{https://arxiv.org/abs/2210.16656}

  \item Lai F., Dai Y., Singapuram S., \textbf{Liu J.}, Zhu X., Madhyastha H., Chowdhury M.
        \textit{FedScale: Benchmarking Model and System Performance of Federated Learning at Scale.}
        ICML 2022. \url{https://arxiv.org/abs/2105.11367}

  \item \textbf{Liu J.*}, Yu P.*, Chowdhury M.
        \textit{Fluid: Resource-Aware Hyperparameter Tuning Engine.}
        MLSys 2021. \url{https://github.com/SymbioticLab/Fluid}

\end{enumerate}
}

% ============================================================
%  PART 2 — NIH BIOGRAPHICAL SKETCH SUPPLEMENT
% ============================================================
\newpage
\fancyhead[L]{\small\textbf{NIH Biographical Sketch Supplement}\enspace OMB No. 0925-0001 (Exp. 12/31/2027)}

{\footnotesize OMB No.~0925-0001 \& 0925-0002 \quad Expiration Dates 12/31/2027 \& 11/30/2027}

\begin{center}
  {\large\bfseries NIH BIOGRAPHICAL SKETCH SUPPLEMENT}\\[2pt]
  {\small (Appended to Biographical Sketch Common Form)}
\end{center}
\vspace{-4pt}

% ── Identifying Information ──────────────────────────────────
\section*{Identifying Information}

\begin{tabular}{@{}L{5.2cm} L{5.0cm} L{4.8cm}@{}}
\textbf{*Name} (Last, First, Middle)  & \textbf{*Persistent Identifier (ORCID)} & \textbf{*Position Title} \\
Liu, Jiachen (Amber)                  & 0000-0003-2317-1956                     & Co-Founder \& CTO \\[6pt]
\multicolumn{3}{@{}l}{\textbf{*Organization and Location}} \\
\multicolumn{3}{@{}l}{Orchestra Research Institute, Inc.\quad|\quad Cambridge, MA, USA}
\end{tabular}

% ── Personal Statement ───────────────────────────────────────
\section*{Personal Statement}
\textit{(Max 3,500 characters. No citations.)}

\vspace{4pt}
I am Co-Founder and Chief Technology Officer of Orchestra Research Institute, an AI-native research platform for coordinating agentic scientific workflows across the full research lifecycle. My role in the DISCOVER project is to contribute evaluation methodology, adversarial threat-surface expertise, and production validation infrastructure grounded in Orchestra's operational deployment of multi-domain AI research agents.

My technical background is directly relevant to this proposal. My doctoral work at the University of Michigan (SymbioticLab, 2020--2025) produced foundational systems for large language model serving, training, and evaluation, as well as open-source tools now widely used in the community. Most relevant to DISCOVER: I co-led the development of Curie, the first AI-agent framework for automated end-to-end scientific experimentation, and EXP-Bench (ICLR 2026), the first peer-reviewed benchmark for evaluating whether AI can autonomously conduct ML research experiments. EXP-Bench provides a principled, reproducible evaluation harness for AI research agents — directly serving as the anchoring benchmark for DISCOVER's Decision Gate Metrics.

As CTO of Orchestra, I oversee the design and operation of multi-domain agentic research pipelines that span literature synthesis, hypothesis generation, experiment execution, and publication. These are precisely the workflows DISCOVER aims to harden against adversarial attack. My experience deploying AI agents in real academic and industry settings — including materials science collaborators at the University of Michigan — gives me direct insight into the threat surfaces and failure modes relevant to Focus Area A. I have not published under a different name.

% ── Honors ───────────────────────────────────────────────────
\section*{Honors}
\textit{(Max 15 entries.)}

\vspace{4pt}
\begin{enumerate}[noitemsep, topsep=2pt, leftmargin=1.5em, label=\arabic*.]
  \item \textbf{ML and Systems Rising Stars} — Program of 2023 (selective national cohort for outstanding Ph.D. students in ML/Systems)
  \item \textbf{Ph.D. Student Fellowship} — University of Michigan, 2020
  \item \textbf{Shanghai Outstanding Graduate (Top 5\%)} — Shanghai Jiao Tong University, 2020
  \item \textbf{Dean's List (Top 5\%)} — Shanghai Jiao Tong University, 2017--2020
  \item \textbf{Undergraduate Scholarship (Top 10\%)} — Shanghai Jiao Tong University
  \item \textbf{University Honors} — University of Michigan, 2019
  \item \textbf{Dean's List} — University of Michigan, 2018 \& 2019
\end{enumerate}

% ── Contributions to Science ─────────────────────────────────
\section*{Contributions to Science}
\textit{(Max 5 contributions; each max 2,000 characters. No citations. No figures or tables.)}
\textit{Products referenced by author/title/year from the Other Significant Products list above.}

\vspace{4pt}

\textbf{1. Benchmarking and Evaluation of AI Research Agents}

The ability to rigorously evaluate AI agents on scientific tasks is a prerequisite for assessing their robustness — including under adversarial conditions. Before EXP-Bench, no peer-reviewed, reproducible benchmark existed for measuring whether AI could autonomously conduct ML research experiments. Evaluation was ad hoc, inconsistent across labs, and not grounded in controlled experimental design.

I co-led the design, implementation, and evaluation of EXP-Bench (Kon \& Liu et al., ICLR 2026), which establishes the first controlled benchmark for this task class. EXP-Bench defines 247 research experiment tasks drawn from published ML papers, with verifiable expected outcomes, isolated execution environments, and reproducibility protocols. My specific contributions included the benchmark architecture, the controlled-environment execution protocol, the task-validity verification pipeline, and the evaluation metrics. EXP-Bench now serves as the anchoring evaluation framework for DISCOVER's Decision Gate Metrics.

I also co-led the Curie framework (Kon \& Liu et al., arXiv 2025), which integrates hypothesis formulation, experiment execution, and result interpretation into a single agentic pipeline with built-in reliability checks. Curie is the production system underlying Orchestra's platform and is the real-world deployment environment for the adversarial evaluation work in this proposal.

\medskip

\textbf{2. Evaluation and Measurement Methodology for AI Systems}

Rigorous system evaluation is central to understanding both capabilities and vulnerabilities. I have built expertise in measurement methodology across multiple AI system contexts: the ML.ENERGY Benchmark (Chung \& Liu et al., NeurIPS 2025 Spotlight) established the first systematic methodology for measuring LLM inference energy consumption at scale, introducing standardized harnesses, reproducible configurations, and the ml.energy public leaderboard. The Evaluation Framework for AI Systems in ``the Wild'' (Jabbour, Liu et al., arXiv 2025) extended evaluation methodology to deployed, real-world AI system conditions — addressing distribution shift, underspecified environments, and behavioral consistency, all of which are directly relevant to adversarial robustness evaluation. My roles in both projects included measurement protocol design, system implementation, and empirical analysis.

\medskip

\textbf{3. ML Systems Infrastructure at Scale}

Robust AI agents for science depend on underlying ML systems that are efficient, reproducible, and well-characterized. My dissertation built a body of work across the ML systems stack: FedScale (Lai \& Liu et al., ICML 2022) is the first large-scale FL benchmark simulating millions of heterogeneous real-world devices, establishing evaluation standards adopted by the FL community. Auxo (Liu et al., SoCC 2023) addresses client selection at scale. Venn (Liu et al., MLSys 2025) provides fair resource management across concurrently running collaborative training jobs. This systems expertise informs my ability to characterize the attack surface of production AI research pipelines and to design defenses that remain practical at operational scale.

% ── Certification ─────────────────────────────────────────────
\section*{Certification}

\vspace{3pt}
{\small\itshape I certify that the information provided is current, accurate, and complete, including information related to domestic and foreign appointments and positions. I also certify that, at the time of submission, I am not a party to a malign foreign talent recruitment program. Misrepresentations and/or omissions may be subject to prosecution and liability pursuant to 18 U.S.C.\ §§\,287, 1001, 1031 and 31 U.S.C.\ §§\,3729--3733 and 3802.}

\vspace{10pt}
\begin{tabular}{@{}L{9cm} L{5.5cm}@{}}
Signature:\enspace\underline{\hspace{7cm}} & Date:\enspace\underline{\hspace{4cm}} \\
\end{tabular}

\vspace{8pt}
{\footnotesize\textit{Note: Signature date must be within 12 months of submission to the Federal funding agency.}}

\end{document}
